\documentclass[12pt,a4paper]{article}

\usepackage[T2A]{fontenc}
\usepackage[utf8]{inputenc}
\usepackage[russian]{babel}
\usepackage{amsmath,amssymb,amsthm,mathtools}
\usepackage{enumitem}
\usepackage{geometry}
\usepackage{hyperref}
\usepackage{stmaryrd}
\usepackage{mathrsfs}

\geometry{margin=2.2cm}

\title{Ответ на экзаменационный вопрос:\\
$\lambda$-исчисление, правила редукции, единственность нормальной формы\\
и правила её достижения, представление рекурсивных функций}
\author{}
\date{}

\newtheorem{definition}{Определение}
\newtheorem{theorem}{Теорема}
\newtheorem{proposition}{Утверждение}
\newtheorem{lemma}{Лемма}
\newtheorem{corollary}{Следствие}
\theoremstyle{remark}
\newtheorem*{remark}{Замечание}
\newtheorem*{example}{Пример}


\title{Билет №13 \\ \small $\lambda$-исчисление, правила редукции, единственность нормальной формы и правила ее достижения, представление рекурсивных функций. (СПбГУ)}
\author{Сериков Владислав Александрович}
\date{28 апреля 2026}


\begin{document}
\maketitle

\section*{Введение}

$\lambda$-исчисление Чёрча --- одна из базовых формальных моделей вычисления.
Оно важно сразу по двум причинам:

\begin{itemize}
    \item с \emph{математической} точки зрения оно даёт строгое определение вычисления;
    \item с \emph{программистской} точки зрения оно лежит в основе функционального программирования.
\end{itemize}

Главная идея проста: вычисление интерпретируется как \emph{редукция выражения}, то есть последовательность локальных преобразований термов.
При этом функции рассматриваются как значения первого класса, а вычисление функции от аргумента реализуется через \emph{подстановку} аргумента в тело функции.

Ниже последовательно изложены:

\begin{enumerate}[label=\arabic*)]
    \item определения $\lambda$-термов;
    \item свободные и связанные переменные;
    \item корректная подстановка;
    \item $\beta$-редукция и стратегии редукции;
    \item нормальные формы, их единственность и достижение;
    \item представление данных и рекурсивных функций в $\lambda$-исчислении.
\end{enumerate}

\section{Синтаксис $\lambda$-исчисления}

\subsection{Определение $\lambda$-термов}

\begin{definition}
Пусть задано счётное множество переменных
\[
x,y,z,u,v,\dots
\]
Множество $\lambda$-термов определяется индуктивно:
\begin{enumerate}
    \item любая переменная $x$ есть $\lambda$-терм;
    \item если $M$ и $N$ --- $\lambda$-термы, то $(MN)$ --- $\lambda$-терм (применение, \emph{application});
    \item если $x$ --- переменная, а $M$ --- $\lambda$-терм, то $(\lambda x.M)$ --- $\lambda$-терм (абстракция, \emph{abstraction}).
\end{enumerate}
\end{definition}

\paragraph{Интуиция.}
\begin{itemize}
    \item $\lambda x.M$ --- функция одного аргумента $x$ с телом $M$;
    \item $MN$ --- применение функции $M$ к аргументу $N$.
\end{itemize}

\subsection{Соглашения о скобках}

Обычно используются следующие сокращения:

\begin{itemize}
    \item применение левоассоциативно:
    \[
    MNP \equiv (MN)P;
    \]
    \item абстракция правоассоциативна:
    \[
    \lambda x.\lambda y.\lambda z.M \equiv \lambda x.(\lambda y.(\lambda z.M));
    \]
    \item область действия символа $\lambda$ максимально велика:
    \[
    \lambda x.MN \equiv \lambda x.(MN).
    \]
\end{itemize}

\paragraph{Минимальные примеры.}
\[
x,\qquad xy,\qquad \lambda x.x,\qquad \lambda x.\lambda y.xy,\qquad (\lambda x.x)y.
\]

\section{Свободные и связанные переменные}

\begin{definition}
Вхождение переменной $x$ в терм называется:
\begin{itemize}
    \item \emph{связанным}, если оно находится в области действия некоторой абстракции $\lambda x$ или само стоит сразу после символа $\lambda$;
    \item \emph{свободным} --- во всех остальных случаях.
\end{itemize}
\end{definition}

\begin{definition}
Множество свободных переменных $FV(M)$ определяется рекурсивно:
\[
FV(x)=\{x\},
\]
\[
FV(MN)=FV(M)\cup FV(N),
\]
\[
FV(\lambda x.M)=FV(M)\setminus\{x\}.
\]
\end{definition}

\begin{definition}
Множество связанных переменных $BV(M)$ определяется рекурсивно:
\[
BV(x)=\varnothing,
\]
\[
BV(MN)=BV(M)\cup BV(N),
\]
\[
BV(\lambda x.M)=BV(M)\cup\{x\}.
\]
\end{definition}

\paragraph{Минимальные примеры.}
\begin{enumerate}
    \item Для терма $x$:
    \[
    FV(x)=\{x\},\qquad BV(x)=\varnothing.
    \]
    \item Для терма $\lambda x.x$:
    \[
    FV(\lambda x.x)=\varnothing,\qquad BV(\lambda x.x)=\{x\}.
    \]
    \item Для терма $x(\lambda x.x)$:
    \[
    FV(x(\lambda x.x))=\{x\},\qquad BV(x(\lambda x.x))=\{x\}.
    \]
\end{enumerate}

\begin{remark}
Одна и та же переменная может иметь и свободные, и связанные вхождения в одном терме.
\end{remark}

\section{Подстановка}

\subsection{Почему наивное определение неверно}

Если понимать подстановку как ``замену всех вхождений $x$ на $N$'', возникают две проблемы:

\begin{enumerate}
    \item \textbf{исчезновение связанной переменной};
    \item \textbf{захват свободной переменной}.
\end{enumerate}

\paragraph{Пример 1: исчезновение связанной переменной.}
Если формально подставить $t$ вместо $x$ в $\lambda x.x$, то получится $\lambda x.t$, а тождественная функция превратится в константу. Это неправильно.

\paragraph{Пример 2: захват свободной переменной.}
Если подставить $y$ вместо $x$ в $\lambda y.x$, то наивно вышло бы $\lambda y.y$. Но свободная переменная $y$ стала связанной, и смысл терма изменился.

Следовательно, подстановка должна учитывать структуру связывания переменных.

\subsection{Корректное определение подстановки}

\begin{definition}
Подстановка $[N/x]M$ определяется рекурсивно:
\begin{align*}
[N/x]x &= N,\\
[N/x]y &= y,\qquad y\neq x,\\
[N/x](PQ) &= ([N/x]P)([N/x]Q),\\
[N/x](\lambda x.P) &= \lambda x.P.
\end{align*}

Если $y\neq x$, то:
\begin{itemize}
    \item если $y\notin FV(N)$, то
    \[
    [N/x](\lambda y.P)=\lambda y.([N/x]P);
    \]
    \item если $y\in FV(N)$, то предварительно выполняется переименование связанной переменной:
    \[
    [N/x](\lambda y.P)=\lambda z.([N/x]([z/y]P)),
    \]
    где $z$ --- свежая переменная, не входящая свободно ни в $N$, ни в $P$.
\end{itemize}
\end{definition}

\subsection{$\alpha$-конверсия}

\begin{definition}
Переименование связанной переменной:
\[
\lambda x.M \rightsquigarrow \lambda y.[y/x]M,
\]
если $y$ не конфликтует со свободными переменными, называется \emph{$\alpha$-конверсией}.
\end{definition}

\begin{definition}
Термы $M$ и $N$ называются \emph{$\alpha$-эквивалентными}, если один получается из другого конечным числом $\alpha$-конверсий. Обозначение:
\[
M\equiv_\alpha N.
\]
\end{definition}

\paragraph{Минимальные примеры.}
\[
\lambda x.x \equiv_\alpha \lambda y.y,
\qquad
\lambda x.\lambda y.x \equiv_\alpha \lambda z.\lambda y.z.
\]

\subsection{Алгоритм корректной подстановки}

\paragraph{Шаги алгоритма подстановки $[N/x]M$.}
\begin{enumerate}
    \item Если $M$ --- переменная:
    \begin{itemize}
        \item если $M=x$, вернуть $N$;
        \item иначе вернуть $M$.
    \end{itemize}
    \item Если $M=P\,Q$, рекурсивно подставить в оба подтерма:
    \[
    [N/x](PQ)=([N/x]P)([N/x]Q).
    \]
    \item Если $M=\lambda x.P$, остановиться:
    \[
    [N/x](\lambda x.P)=\lambda x.P,
    \]
    потому что внутри этой абстракции $x$ уже связан.
    \item Если $M=\lambda y.P$ и $y\neq x$:
    \begin{itemize}
        \item если $y\notin FV(N)$, подставить в тело:
        \[
        [N/x](\lambda y.P)=\lambda y.([N/x]P);
        \]
        \item если $y\in FV(N)$, выбрать свежую переменную $z$, выполнить $\alpha$-переименование
        \[
        \lambda y.P \equiv_\alpha \lambda z.([z/y]P),
        \]
        и затем продолжить подстановку.
    \end{itemize}
\end{enumerate}

\paragraph{Иллюстрация.}
Вычислим:
\[
[y/x](\lambda y.x).
\]
Поскольку $y\in FV(y)$, возникает захват. Поэтому сначала переименуем:
\[
\lambda y.x \equiv_\alpha \lambda z.x.
\]
Теперь
\[
[y/x](\lambda z.x)=\lambda z.[y/x]x=\lambda z.y.
\]

\subsection{Корректность подстановки}

\begin{proposition}
Корректная подстановка не захватывает свободные переменные подставляемого терма.
\end{proposition}

\begin{proof}
Доказывается по индукции по структуре терма $M$.

\textbf{База.}
\begin{itemize}
    \item Если $M=x$, то $[N/x]x=N$; захвата нет, так как новых связывающих абстракций не появляется.
    \item Если $M=y\neq x$, то $[N/x]y=y$; ничего не меняется.
\end{itemize}

\textbf{Индукционный переход.}
\begin{itemize}
    \item Для применения $M=PQ$ подстановка выполняется независимо в $P$ и $Q$. По предположению индукции в каждом из них захвата не возникает, следовательно, не возникает и во всём $PQ$.
    \item Для абстракции $M=\lambda x.P$ подстановка внутрь не идёт; значит, захват невозможен.
    \item Для абстракции $M=\lambda y.P$, $y\neq x$:
    \begin{itemize}
        \item если $y\notin FV(N)$, то связывающая переменная $y$ не может захватить свободные переменные терма $N$;
        \item если $y\in FV(N)$, то перед подстановкой выполняется $\alpha$-переименование в свежую переменную $z$, которая не входит свободно ни в $N$, ни в $P$. Следовательно, после переименования захват невозможен.
    \end{itemize}
\end{itemize}
Тем самым подстановка определена корректно.
\end{proof}

\section{$\beta$-редукция и вычисление}

\subsection{Редексы}

\begin{definition}
Терм вида
\[
(\lambda x.M)N
\]
называется \emph{$\beta$-редексом} или \emph{редуцируемым выражением}.
\end{definition}

Смысл: это явное применение функции $\lambda x.M$ к аргументу $N$.

\subsection{$\beta$-редукция}

\begin{definition}
\emph{$\beta$-редукцией} называется преобразование
\[
(\lambda x.M)N \;\to_\beta\; [N/x]M.
\]
\end{definition}

\paragraph{Минимальные примеры.}
\[
(\lambda x.x)y \to_\beta y,
\]
\[
(\lambda x.y)t \to_\beta y,
\]
\[
(\lambda x.x(xy))N \to_\beta N(Ny).
\]

\subsection{Алгоритм одного шага $\beta$-редукции}

\paragraph{Шаги.}
\begin{enumerate}
    \item Найти в терме подвыражение вида $(\lambda x.M)N$.
    \item Убедиться, что подстановка $[N/x]M$ выполняется корректно (при необходимости сделать $\alpha$-переименование).
    \item Заменить редекс результатом подстановки.
\end{enumerate}

\paragraph{Иллюстрация.}
\[
(\lambda x.\lambda y.yx)zv
\]
Сначала редуцируем самый левый внешний редекс:
\[
(\lambda x.\lambda y.yx)zv \to_\beta (\lambda y.yz)v \to_\beta vz.
\]

\section{Нормальные формы}

\begin{definition}
$\lambda$-терм находится в \emph{нормальной форме}, если не содержит ни одного $\beta$-редекса.
\end{definition}

\paragraph{Примеры.}
\begin{itemize}
    \item $x$, $\lambda x.x$, $\lambda x.\lambda y.xy$ --- нормальные формы;
    \item $(\lambda x.x)y$ --- не нормальная форма.
\end{itemize}

\subsection{Слабая и полная нормальная форма}

\begin{definition}
Терм находится в \emph{слабой (головной) нормальной форме}, если у него нет \emph{внешних} редексов.
\end{definition}

Это важно, потому что некоторые стратегии не редуцируют внутри абстракций.

\section{Стратегии редукции}

Поскольку в одном терме может быть несколько редексов, надо уточнить, \emph{какой} редекс редуцировать.

\subsection{Нормальный порядок}

\begin{definition}
\emph{Нормальный порядок редукции} --- это стратегия, при которой на каждом шаге выбирается \emph{самый левый внешний} редекс.
\end{definition}

\paragraph{Что значит ``самый левый внешний''.}
\begin{itemize}
    \item \emph{внешний} --- не содержащийся внутри другого редекса;
    \item \emph{левый} --- первый слева среди внешних.
\end{itemize}

\subsection{Вызов по имени}

\begin{definition}
\emph{Вызов по имени} (\emph{call by name}) --- как нормальный порядок, но без редукции внутри абстракций.
\end{definition}

\subsection{Вызов по необходимости}

\begin{definition}
\emph{Вызов по необходимости} (\emph{call by need}) --- ленивый вариант вызова по имени с мемоизацией: если аргумент вычислен один раз, затем используется уже вычисленное значение.
\end{definition}

\subsection{Вызов по значению}

\begin{definition}
\emph{Вызов по значению} (\emph{call by value}) --- сначала вычисляются аргументы, затем применяется функция.
\end{definition}

Это \emph{строгая} стратегия: аргумент вычисляется даже тогда, когда в теле функции он не используется.

\subsection{Сравнение стратегий на минимальном примере}

Рассмотрим
\[
(\lambda y.v)\big((\lambda x.xxy)(\lambda x.xxy)\big).
\]

\paragraph{Нормальный порядок.}
Сначала редуцируется внешний редекс:
\[
(\lambda y.v)\Omega \to_\beta v,
\]
где
\[
\Omega=((\lambda x.xxy)(\lambda x.xxy)).
\]
Результат получен сразу.

\paragraph{Аппликативный порядок (вызов по значению).}
Сначала пытаемся редуцировать аргумент $\Omega$:
\[
\Omega \to_\beta \Omega y \to_\beta \Omega yy \to_\beta \cdots
\]
Получаем бесконечное вычисление.

\paragraph{Вывод.}
Неправильная стратегия может не привести к нормальной форме даже тогда, когда нормальная форма существует.

\section{Единственность нормальной формы}

\subsection{Проблема}

Если редексов несколько, можно выбирать их по-разному. Возникают вопросы:

\begin{enumerate}
    \item Всегда ли разные последовательности редукций приводят к одному и тому же результату?
    \item Если нормальная форма существует, можно ли её гарантированно достичь?
\end{enumerate}

Ключ к ответу --- теорема Чёрча--Россера.

\subsection{Теорема Чёрча--Россера}

\begin{theorem}[Чёрча--Россера]
Если
\[
P \to_\beta^* N
\qquad\text{и}\qquad
P \to_\beta^* M,
\]
то существует такой терм $T$, что
\[
N \to_\beta^* T
\qquad\text{и}\qquad
M \to_\beta^* T.
\]
\end{theorem}

Это свойство также называют \emph{свойством ромба} или \emph{конфлюэнтностью}.

\subsection{Следствие: единственность нормальной формы}

\begin{corollary}
Если терм имеет нормальную форму, то она единственна с точностью до $\alpha$-эквивалентности.
\end{corollary}

\begin{proof}
Пусть терм $P$ имеет две нормальные формы $N$ и $M$, то есть
\[
P \to_\beta^* N,
\qquad
P \to_\beta^* M,
\]
причём $N$ и $M$ редексов не содержат.

По теореме Чёрча--Россера существует терм $T$, такой что
\[
N \to_\beta^* T,
\qquad
M \to_\beta^* T.
\]
Но $N$ --- нормальная форма, значит, из неё уже нельзя сделать ни одного $\beta$-шага. Следовательно, единственно возможна пустая цепочка редукций, то есть
\[
N \equiv_\alpha T.
\]
Аналогично
\[
M \equiv_\alpha T.
\]
Отсюда
\[
N \equiv_\alpha M.
\]
Следовательно, нормальная форма единственна с точностью до переименования связанных переменных.
\end{proof}

\subsection{Почему это доказательство корректно}

Оно использует два факта:

\begin{enumerate}
    \item теорема Чёрча--Россера гарантирует \emph{сходимость} любых двух путей редукции;
    \item нормальная форма \emph{не редуцируется дальше}.
\end{enumerate}

Значит, если два пути привели к двум нормальным формам, они обязаны совпасть с точностью до $\alpha$-конверсии.

\section{Правила достижения нормальной формы}

\subsection{Главная теорема о нормальном порядке}

\begin{theorem}
Если терм имеет нормальную форму, то нормальный порядок редукции обязательно приведёт к ней.
\end{theorem}

Это один из центральных результатов: \emph{нормальный порядок нормализует}, если нормализация вообще возможна.

\subsection{Смысл теоремы}

\begin{itemize}
    \item Не всякая стратегия годится.
    \item Но стратегия ``левый внешний редекс'' годится всегда.
\end{itemize}

Именно поэтому нормальный порядок считается \emph{правилом достижения нормальной формы}.

\subsection{Практический алгоритм приведения к нормальной форме}

\paragraph{Алгоритм.}
Для терма $M$:
\begin{enumerate}
    \item Найти самый левый внешний редекс.
    \item Выполнить один $\beta$-шаг.
    \item Повторять, пока редексов не останется.
\end{enumerate}

\paragraph{Корректность алгоритма.}
\begin{itemize}
    \item Если алгоритм завершился, полученный терм не содержит редексов, то есть является нормальной формой.
    \item Если у исходного терма нормальная форма существовала, теорема о нормальном порядке гарантирует, что алгоритм завершится именно на ней.
    \item По единственности нормальной формы результат не зависит от других возможных успешных путей редукции.
\end{itemize}

\paragraph{Следовательно:}
нормальный порядок --- корректное и полное правило достижения нормальной формы.

\subsection{Минимальная иллюстрация}

Рассмотрим терм
\[
(\lambda x.(\lambda y.yx)z)v.
\]

\paragraph{Шаг 1.}
Самый левый внешний редекс:
\[
(\lambda x.(\lambda y.yx)z)v \to_\beta (\lambda y.yv)z.
\]

\paragraph{Шаг 2.}
Следующий редекс:
\[
(\lambda y.yv)z \to_\beta zv.
\]

\paragraph{Шаг 3.}
Терм $zv$ редексов не содержит, это нормальная форма.

\section{Термы без нормальной формы}

Не всякий терм нормализуем.

\paragraph{Классический пример.}
\[
\Omega=(\lambda x.xx)(\lambda x.xx).
\]
Тогда
\[
\Omega \to_\beta \Omega.
\]
Значит, редукция бесконечна, и нормальной формы нет.

\paragraph{Вывод.}
Задача достижения нормальной формы имеет смысл только для тех термов, которые \emph{нормализуемы}.

\section{Представление данных в $\lambda$-исчислении}

\subsection{Булевы значения}

Определим:
\[
\mathsf{true}=\lambda x.\lambda y.x,
\qquad
\mathsf{false}=\lambda x.\lambda y.y.
\]

\paragraph{Интерпретация.}
\begin{itemize}
    \item $\mathsf{true}$ возвращает первый аргумент;
    \item $\mathsf{false}$ возвращает второй аргумент.
\end{itemize}

\subsection{Условный оператор}

\[
\mathsf{if}\ C\ \mathsf{then}\ E_1\ \mathsf{else}\ E_2
\;:=\;
C\,E_1\,E_2.
\]

\paragraph{Проверка.}
\[
\mathsf{true}\ E_1\ E_2
=
(\lambda x.\lambda y.x)E_1E_2
\to_\beta E_1,
\]
\[
\mathsf{false}\ E_1\ E_2
=
(\lambda x.\lambda y.y)E_1E_2
\to_\beta E_2.
\]

\subsection{Логические операции}

\[
\mathsf{not}\ P := P\ \mathsf{false}\ \mathsf{true},
\]
\[
P\ \mathsf{and}\ Q := P\ Q\ \mathsf{false},
\]
\[
P\ \mathsf{or}\ Q := P\ \mathsf{true}\ Q.
\]

\subsection{Пары}

\[
(E_1,E_2):=\lambda z.zE_1E_2.
\]

\[
\mathsf{fst}\ P := P\ \mathsf{true},
\qquad
\mathsf{snd}\ P := P\ \mathsf{false}.
\]

\paragraph{Проверка.}
\[
\mathsf{fst}(E_1,E_2)
=
(\lambda z.zE_1E_2)\mathsf{true}
\to_\beta
\mathsf{true}E_1E_2
\to_\beta E_1.
\]
Аналогично
\[
\mathsf{snd}(E_1,E_2)\to_\beta E_2.
\]

\section{Нумералы Чёрча}

\subsection{Определение}

Натуральное число $n$ кодируется термом
\[
\overline{n}:=\lambda f.\lambda x.f^n x,
\]
то есть
\[
\overline{0}=\lambda f.\lambda x.x,
\]
\[
\overline{1}=\lambda f.\lambda x.fx,
\]
\[
\overline{2}=\lambda f.\lambda x.f(fx),
\]
\[
\overline{3}=\lambda f.\lambda x.f(f(fx)),
\]
и т.д.

\subsection{Интуиция}

Число $\overline{n}$ означает: \emph{применить функцию $f$ к $x$ ровно $n$ раз}.

\paragraph{Минимальный пример.}
\[
\overline{2}ft
=
(\lambda f.\lambda x.f(fx))ft
\to_\beta
f(ft).
\]

\section{Арифметика в $\lambda$-исчислении}

\subsection{Функция следования}

\[
\mathsf{succ}
=
\lambda n.\lambda f.\lambda x.f(nfx).
\]

\paragraph{Проверка на примере.}
\[
\mathsf{succ}\ \overline{2}
=
\lambda f.\lambda x.f(\overline{2}fx)
\to_\beta
\lambda f.\lambda x.f(f(fx))
=
\overline{3}.
\]

\subsection{Сложение}

\[
\mathsf{add}
=
\lambda m.\lambda n.\lambda f.\lambda x.mf(nfx).
\]

\paragraph{Интуиция.}
Сначала $n$ раз применяем $f$ к $x$, затем ещё $m$ раз.

\paragraph{Минимальный пример.}
\[
\mathsf{add}\ \overline{1}\ \overline{2}
\to_\beta
\overline{3}.
\]

\subsection{Умножение}

\[
\mathsf{mul}
=
\lambda m.\lambda n.\lambda f.\lambda x.m(nf)x.
\]

\paragraph{Интуиция.}
Функция $nf$ означает ``применить $f$ $n$ раз'', а затем $m$ раз применить этот блок.

\section{Проверка нуля}

\[
\mathsf{isZero}
=
\lambda n.n(\lambda x.\mathsf{false})\mathsf{true}.
\]

\paragraph{Корректность.}
\begin{itemize}
    \item Для $\overline{0}$:
    \[
    \mathsf{isZero}\ \overline{0}
    =
    (\lambda n.n(\lambda x.\mathsf{false})\mathsf{true})\overline{0}
    \to_\beta
    \mathsf{true}.
    \]
    \item Для любого $\overline{n}$, $n>0$, функция $(\lambda x.\mathsf{false})$ применяется хотя бы один раз, и результатом будет $\mathsf{false}$.
\end{itemize}

\section{Предшественник и вычитание}

\subsection{Идея Клини}

Предшественник кодируется существенно сложнее, чем следование. Используется вспомогательная функция
\[
\mathsf{step}(m,n)=(n,n+1),
\]
которая по паре строит следующую.

Если начать с пары $(0,0)$ и применить $\mathsf{step}$ $m$ раз, получим
\[
(0,0)\mapsto (0,1)\mapsto (1,2)\mapsto (2,3)\mapsto \cdots \mapsto (m-1,m).
\]
Тогда первый компонент и будет предшественником.

\subsection{Определения}

\[
\mathsf{step}
=
\lambda p.(\mathsf{snd}\ p,\ \mathsf{succ}(\mathsf{snd}\ p)).
\]

\[
\mathsf{pred}
=
\lambda m.\mathsf{fst}(m\ \mathsf{step}\ (0,0)).
\]

\[
\mathsf{sub}
=
\lambda m.\lambda n.n\ \mathsf{pred}\ m.
\]

\paragraph{Интуиция.}
Чтобы вычесть $n$ из $m$, достаточно $n$ раз применить $\mathsf{pred}$ к $m$.

\section{Представление рекурсивных функций}

\subsection{Проблема рекурсии}

В обычном языке программирования можно написать:
\[
\mathsf{fact}(n)=
\begin{cases}
1,& n=0,\\
n\cdot \mathsf{fact}(n-1),& n>0.
\end{cases}
\]
Но в чистом $\lambda$-исчислении у функций \emph{нет имён}, поэтому нельзя просто сослаться на $\mathsf{fact}$ внутри её же тела.

\subsection{Комбинатор неподвижной точки}

\begin{definition}
Терм $Y$ называется \emph{комбинатором неподвижной точки}, если для любого $f$
\[
Yf =_\beta f(Yf).
\]
\end{definition}

Смысл: $Yf$ --- неподвижная точка функции $f$.

\subsection{Комбинатор Тьюринга}

Положим
\[
A=\lambda x.\lambda y.y(xxy),
\qquad
Y_T=AA.
\]

\begin{proposition}
$Y_T$ является комбинатором неподвижной точки.
\end{proposition}

\begin{proof}
\[
Y_Tf
=
AAf
=
(\lambda x.\lambda y.y(xxy))Af
\to_\beta
(\lambda y.y(AAy))f
\to_\beta
f(AAf)
=
f(Y_Tf).
\]
\end{proof}

\subsection{Как из этого получается рекурсия}

Пусть
\[
H=
\lambda f.\lambda n.\mathsf{if}\ (\mathsf{isZero}\ n)\ \mathsf{then}\ \overline{1}\ \mathsf{else}\ \mathsf{mul}\ n\ (f(\mathsf{pred}\ n)).
\]

Тогда факториал определяется как
\[
\mathsf{fact}=Y_T H.
\]

Действительно,
\[
\mathsf{fact}
=
Y_T H
=_\beta
H(Y_T H)
=
H(\mathsf{fact}),
\]
то есть
\[
\mathsf{fact}(n)
=
\mathsf{if}\ (\mathsf{isZero}\ n)\ \mathsf{then}\ \overline{1}\ \mathsf{else}\ \mathsf{mul}\ n\ (\mathsf{fact}(\mathsf{pred}\ n)).
\]

Это в точности рекурсивное определение факториала.

\subsection{Алгоритм представления рекурсивной функции}

\paragraph{Шаги.}
Пусть требуется определить рекурсивную функцию
\[
F(x)=\Phi(F,x).
\]
Тогда:

\begin{enumerate}
    \item Выделяем \emph{шаговый функционал}
    \[
    H=\lambda f.\lambda x.\Phi(f,x),
    \]
    где рекурсивный вызов заменён параметром $f$.
    \item Применяем комбинатор неподвижной точки:
    \[
    F=YH.
    \]
    \item Проверяем:
    \[
    F=YH=_\beta H(YH)=H(F),
    \]
    то есть рекурсивное уравнение выполнено.
\end{enumerate}

\subsection{Минимальный пример: сумма чисел от $1$ до $n$}

Определим
\[
H_{\Sigma}
=
\lambda f.\lambda n.
\mathsf{if}\ (\mathsf{isZero}\ n)\ \mathsf{then}\ \overline{0}\ \mathsf{else}\ \mathsf{add}\ n\bigl(f(\mathsf{pred}\ n)\bigr).
\]
Тогда
\[
\mathsf{sum}=YH_{\Sigma}.
\]
Получаем
\[
\mathsf{sum}(n)=
\begin{cases}
0,& n=0,\\
n+\mathsf{sum}(n-1),& n>0.
\end{cases}
\]

\section{Корректность представления рекурсивных функций}

\begin{proposition}
Если $Y$ --- комбинатор неподвижной точки, то терм $F:=YH$ является решением функционального уравнения
\[
F=HF.
\]
\end{proposition}

\begin{proof}
По определению комбинатора неподвижной точки
\[
YH =_\beta H(YH).
\]
Положив $F:=YH$, получаем
\[
F =_\beta HF.
\]
То есть $F$ действительно является фиксированной точкой функционала $H$.
\end{proof}

\begin{proposition}
Если функционал $H$ построен из уже корректно представленных операций над числами, булевыми значениями и парами, то $YH$ корректно представляет соответствующую рекурсивную функцию.
\end{proposition}

\begin{proof}
Индукция по схеме определения функции.

\textbf{База.}
Элементарные операции ($\mathsf{isZero}$, $\mathsf{succ}$, $\mathsf{add}$, $\mathsf{mul}$, $\mathsf{pred}$ и т.д.) уже заданы термами, корректность которых проверяется непосредственной $\beta$-редукцией.

\textbf{Переход.}
Пусть функция определяется рекурсивно как фиксированная точка функционала $H$. Тогда по предыдущему утверждению терм $F=YH$ удовлетворяет уравнению $F=_\beta HF$. Но правая часть $HF$ реализует ровно тот рекурсивный шаг, который был задан в определении функции. Следовательно, $F$ вычисляет ту же самую функцию.
\end{proof}

\section{Связь с частично рекурсивными функциями}

В теории вычислимости под \emph{рекурсивными функциями} обычно понимают вычислимые числовые функции, строящиеся из базовых функций с помощью композиции, примитивной рекурсии и минимизации.

В рамках $\lambda$-исчисления эти функции представимы потому, что:

\begin{enumerate}
    \item данные (числа, булевы значения, пары) кодируются термами;
    \item элементарные арифметические и логические операции выражаются $\lambda$-термами;
    \item рекурсия реализуется через комбинаторы неподвижной точки.
\end{enumerate}

Следовательно, $\lambda$-исчисление способно выражать общий рекурсивный процесс и потому является полноценной моделью вычислений.

\section*{Краткий конспективный итог}

\begin{enumerate}
    \item $\lambda$-термы строятся из переменных, применения и абстракции.
    \item Вычисление в $\lambda$-исчислении --- это $\beta$-редукция:
    \[
    (\lambda x.M)N \to_\beta [N/x]M.
    \]
    \item Подстановка должна быть \emph{корректной}, то есть не допускать захвата свободных переменных; для этого используется $\alpha$-конверсия.
    \item Нормальная форма --- терм без редексов.
    \item По теореме Чёрча--Россера нормальная форма, если существует, единственна с точностью до $\alpha$-эквивалентности.
    \item Нормальный порядок редукции (левый внешний редекс) всегда достигает нормальной формы, если она существует.
    \item Булевы значения, пары и натуральные числа кодируются внутри самого $\lambda$-исчисления.
    \item Рекурсивные функции представляются с помощью комбинаторов неподвижной точки:
    \[
    F=YH.
    \]
\end{enumerate}

\section*{Минимальный набор формул для запоминания}

\[
(\lambda x.M)N \to_\beta [N/x]M
\]

\[
FV(x)=\{x\},\qquad
FV(MN)=FV(M)\cup FV(N),\qquad
FV(\lambda x.M)=FV(M)\setminus\{x\}
\]

\[
\mathsf{true}=\lambda x.\lambda y.x,
\qquad
\mathsf{false}=\lambda x.\lambda y.y
\]

\[
\overline{0}=\lambda f.\lambda x.x,
\qquad
\overline{n}=\lambda f.\lambda x.f^n x
\]

\[
\mathsf{succ}=\lambda n.\lambda f.\lambda x.f(nfx)
\]

\[
\mathsf{add}=\lambda m.\lambda n.\lambda f.\lambda x.mf(nfx)
\]

\[
\mathsf{mul}=\lambda m.\lambda n.\lambda f.\lambda x.m(nf)x
\]

\[
\mathsf{isZero}=\lambda n.n(\lambda x.\mathsf{false})\mathsf{true}
\]

\[
Yf =_\beta f(Yf)
\]

\[
\mathsf{fact}=Y\Bigl(\lambda f.\lambda n.\mathsf{if}\ (\mathsf{isZero}\ n)\ \mathsf{then}\ \overline{1}\ \mathsf{else}\ \mathsf{mul}\ n\bigl(f(\mathsf{pred}\ n)\bigr)\Bigr)
\]

\end{document}
